\input{preamble}

\begin{document}

\title{Basic math}
\maketitle

\phantomsection
\label{section-phantom}

\tableofcontents

\section{Signature}
\label{section-signature}

The signature of a form quadratic form can be
equivalently computed as the
maximal dimension of an isotropic subspace,
or as the maximal dimension of the
subspace where the form is negative.

\begin{remark}
\label{remark-nondegenerate-vs-anisotropic}
Let $B:V \times V \to k$ be a bilinear form.
The {\it radical} of $B$ is
$$
\operatorname{rad}(B)=
\{v \in V : B(v,w)=0\text{ for all }w
\in V\}.
$$
The form is {\it nondegenerate} if
$\operatorname{rad}(B)=0$. In coordinates,
if $A$ is the Gram matrix of $B$, this is
equivalent to $\det(A)\neq 0$.

This should not be confused with the
condition that there are no nonzero isotropic
vectors. A vector $v\neq 0$ is {\it
isotropic} if $B(v,v)=0$, and a form with no
nonzero isotropic vectors is called
{\it anisotropic}. Anisotropic implies
nondegenerate, since a nonzero vector in the
radical would be isotropic. The converse is
false; for example the form with matrix
$$
\begin{pmatrix}
1&0\\
0&-1
\end{pmatrix}
$$
is nondegenerate, but $(1,1)$ is isotropic.
\end{remark}

\section{Adjoint}
\label{section-adjoint}

\begin{definition}
\label{definition-adjoint}
Let  $A:(V,\left\langle{}\cdot,\cdot\right\rangle{}_V) \to
(W,\left\langle{}\cdot,\cdot\right\rangle{}_W)$ be a
linear transformation of inner product vector
spaces, its {\it adjoint} is the
map
$$
A^{\dag}:W\to V
$$
satisfying
\begin{equation}
\label{equation-adjoint}
\left\langle{}Av,w\right\rangle{}_W=\left\langle{}v,A^\dag w\right\rangle{}_V
\end{equation}
\end{definition}

Let $G$ be a matrix representation for
$\left\langle{}\cdot,\cdot\right\rangle{}_V$ and $H$ for
$\left\langle{}\cdot,\cdot\right\rangle{}_W$. Then Eq.
\ref{equation-adjoint} is
$$
v^{\mathbf{T}}A^{\mathbf{T}}Hw=(Av)
^{\mathbf{T}}Hw=v^{\mathbf{T}}GA^\dag
w
$$
So
\begin{equation}
\label{equation-adjoint-transpose}
A^{\mathbf{T}}H=GA^\dag \iff
A^\dag:=G^{-1}A^{\mathbf{T}}H
\end{equation}
which gives the relationship between the
usual transpose matrix and the adjoint.

\begin{definition}
\label{definition-Frobenius-norm}
In the space of matrices we define {\it
Frobenius norm} as
\begin{equation}
\label{equation-Frobenius-norm}
\|A\|^2=\text{tr}A^\dag
A=\text{tr}(G^{-1}A^{\mathbf{T}}HA
\end{equation}
\end{definition}

\section{Fréchet derivative}
\label{section-Fréchet-derivative}

\begin{definition}
\label{definition-Frechet-derivative}
\cite{zo1}, Section 8.2.1, Definition 1. A
function $f:E \to \mathbb{R}^n$
defined on a set $E \subset \mathcal{R}^m$ is
{\it differentiable} at the point
$x \in E$, which is a limit point of $E$, if
$$
f(x+h)-f(x)=L(x)h+\alpha(x;h)
$$
where $L(x):\mathbb{R}^m\to \mathbb{R}^n$ is
a function that is linear in $h$
and $\alpha(x;h)=o(h)$ as $h\to 0$, $x+h\in
E$.
\end{definition}

This means that
$$
\lim_{\|h\|\to 0}
\frac{\|f(x+h)-f(x)-Ah\|}{\|h\|}=0
$$
\begin{lemma}
\label{lemma-differentiating-bilinear-form}
Let $V$ be a finite dimensional real vector
space and let $B:V\times V\to \mathbb{R}$ be
a bilinear form. If $a(t),b(t)$ are
differentiable
curves in $V$, then
$$
\frac{d}{dt}B(a(t),b(t))
=
B(a'(t),b(t))+B(a(t),b'(t)).
$$
\end{lemma}

\begin{proof}
Choose a basis $e_1,\ldots,e_n$ of $V$. Write
$$
a(t)=\sum_i a_i(t)e_i,\qquad
b(t)=\sum_j b_j(t)e_j.
$$
If $B_{ij}=B(e_i,e_j)$, then by bilinearity
$$
B(a(t),b(t))=\sum_{i,j}B_{ij}a_i(t)b_j(t).
$$
Therefore the ordinary product rule gives
$$
\frac{d}{dt}B(a(t),b(t))
=
\sum_{i,j}B_{ij}a_i'(t)b_j(t)
+
\sum_{i,j}B_{ij}a_i(t)b_j'(t).
$$
Packing the two sums back into
coordinate-free notation, this is
$$
B(a'(t),b(t))+B(a(t),b'(t)).
$$
\end{proof}

\section{Arzelà-Ascoli Theorem}
\label{section-Arzela-Ascoli-theorem}
\begin{definition}[Equicontinuous family]
\label{definition-equicontinuous}
From Wikipedia. Let $X$ be a compact
Hausdorff space and let  $C(X)$ be the
space of real-valued continuous functions on
$X$. A subset $\mathcal{F} \subset
C(X)$ is said to be {\it equicontinuous} if
for every $x in X$ and every
$\varepsilon>0$, $x$ has a neighbourhood
$U_x$ such that
$$
\forall y \in U_x, \forall f \in \mathcal{F}:
\qquad |f(y)-f(x)|<\varepsilon
$$
\end{definition}

\begin{definition}[Pointwise bounded]
\label{definition-pointwise-bounded}
A set $\mathcal{F}\subset C(X,\mathbb{R})$ is
said to be {\it pointwise
bounded} if for every $x \in X$,
$$
\text{sup}\{|f(x)|:f \in \mathcal{F}\}<\infty
$$
\end{definition}

\begin{theorem}[Arzelà-Ascoli]
\label{theorem-Arzela-Ascoli}
Let $X$ be a compact Hausdorff space. Then a
subset $\mathcal{F}$ of $C(X)$ is
relatively compact in the topology induced by
the uniform norm if and only if it
is equicontinuous and pointwise bounded.
\end{theorem}

\bibliography{my}
\bibliographystyle{amsalpha}

\end{document}
